\section{Conclusion}
\label{sec:conclusion}

We have described a protocol that makes classical a posteriori eigenvalue
bounds computable for neural eigenpairs of the discrete Schr\"odinger
operator, and an implementation of it that runs on ordinary hardware. The
essential element is not any of the inequalities---they are textbook---but the
verification of their hypothesis: by counting the spectrum below two shifts
with Sturm sequences or Sylvester inertia, the interval that the theorems
require is constructed and checked rather than assumed. The same machinery
supplies the spectral gap that turns a residual into a two-sided eigenvalue
enclosure and an eigenfunction angle bound. At a degenerate level the same
machinery is applied to the invariant subspace rather than to a state: the
multiplicity is measured by the count, the cluster receives a two-sided
enclosure, and the angle between the trial subspace and the true one is
bounded, so that the structural obstacle of degeneracy becomes a certificate
rather than a refusal.

The experiments on six potentials in one and two dimensions give a concrete
picture of what can and cannot be certified. All 39 filtered enclosures we
report contain the discrete eigenvalue they refer to---18 from the main sweep
and 21 from the leakage audit, the two-dimensional ones checked against a
dense LAPACK solve rather than the ARPACK values that propose the
intervals---and the certified width is
between $5$ and $40$ times the true error on the harmonic and anharmonic
potentials, between $50$ and $3\times10^{2}$ times on the infinite well and
$7$--$8\times10^{2}$ on the box---where width and error are both at
round-off---and up to $3\times10^{3}$ times on the finite well, where the
bound is correct but uninformative. A
degenerate state with a $4\times10^{-12}$ residual is refused as a state while
the same eigenspace is accepted as a two-dimensional invariant subspace, with a
measured multiplicity of two and a principal-angle bound of
$8\times10^{-13}$ on exact inputs and $6\times10^{-3}$ on network inputs. Of
four independently trained networks on that level, three collapsed onto the
same representative and one found its partner; the certificate refused the
pairs that do not span the level and accepted the pairs that do. Splitting the
level by a controlled anisotropy locates the boundary between the two
certificates exactly at the counting resolution---a split of $10^{-6}$ is
refused as a state just as a split of $10^{-13}$ is, while the subspace is
certified in both cases---so the obstacle is computable rather than
structural. The network
error on the $81\times81$ grid is smaller than the discretization gap, so the
certificate correctly reports the grid as the limiting factor.

The implementation, the drivers that regenerate every table, and the test
suite that locks the guards are available with this article. The two open
items of the previous version of this work are closed here in part. The count
is now computed in directed-rounding interval arithmetic---provably in one
dimension, and in two dimensions whenever the potential is separable, which
covers two of our three two-dimensional problems; the general sparse
factorization still refuses too often in intervals to be useful, and we
report the measurement rather than hide it. And the continuous operator now
receives an upper bound from the variational principle---exact in exact
arithmetic, and limited only by the floating-point evaluation of the
quotient---at the price of a Rayleigh quotient rather than of an interpolation
constant; the lower
side exists only where a comparison principle does, which we demonstrate for
the infinite well and the box and show to fail for the finite well. Neither addition changes the
protocol, and neither touches the cluster construction: what limits the
certificate at a degenerate level is the resolution at which two shifts can
be separated, not the requirement that the level be simple.
